Los conjuntos de datos son vitales a la hora de desarrollar y evaluar IDS basados en aprendizaje automático, ya que contienen las muestras a partir de las cuales los modelos aprenden los patrones para clasificar correctamente. Un conjunto de datos que sirva como base experimental debe incluir tráfico benigno y malicioso etiquetado correctamente y contar con una amplia variedad de muestras. También es importante que cuente con características útiles y descriptivas y que presente flujos o paquetes realistas \cite{sharafaldin2018cicids2017}.

Uno de los conjuntos de datos más utilizados a lo largo de la historia es el KDD Cup 99. Fue creado a partir de los datos generados por el programa DARPA 1998 y las muestras están descritas por 41 características. Cada muestra está etiquetada como tráfico normal o como un ataque específico. Estos ataques específicos se dividen en denegación de servicio, ataques de usuario a root, ataques remotos a local y ataques de reconocimiento. Al ser un dataset público, ha sido utilizado durante años para evaluar IDS basados en aprendizaje automático \cite{tavallaee2009kdd,buczak2016survey}. No obstante, cuenta con varias limitaciones. A día de hoy, han pasado varias décadas desde su creación, por lo que no contiene ataques actuales ni características modernas. También contiene un gran número de registros redundantes, introduciendo la posibilidad de que a la hora de entrenar el modelo se produzca un sesgamiento hacia las clases mayoritarias. En este contexto surge el NSL-KDD, el cual es un conjunto derivado de KDD Cup 99 que elimina sus registros duplicados y reduce el tamaño del conjunto. No obstante, sigue heredando algunas limitaciones del conjunto original, destacando su antigüedad \cite{tavallaee2009kdd,moustafa2015unswnb15}.

Con el objetivo de obtener un conjunto de datos que refleje escenarios modernos, surge UNSW-NB15. Este fue creado mezclando tráfico real moderno con tráfico generado artificialmente. Para crear este tráfico artificial se utilizó la herramienta IXIA PerfectStorm. El tráfico real fue capturado en formato PCAP mediante la herramienta tcpdump y posteriormente procesado con las herramientas Argus, Bro-IDS y algoritmos propios para añadir características adicionales. Este conjunto de datos contiene tráfico benigno y nueve tipos distintos de ataques. En comparación con KDD Cup 99 y NSL-KDD, este conjunto de datos ofrece una alternativa más actual para IDS contemporáneos \cite{moustafa2015unswnb15}.

A día de hoy, uno de los más relevantes es el CIC-IDS2017. Este es un conjunto de datos público desarrollado por el Canadian Institute for Cybersecurity, el cual proporciona datos completos y actualizados. Este fue diseñado teniendo en cuenta diversos factores, como la diversidad de ataques, necesidad de una captura completa, etiquetado correcto, disponibilidad de protocolos, metadatos y proporcionar un conjunto de características sólido y descriptivo. Incluye tráfico benigno y distintos tipos de ataques, entre los que se pueden encontrar DoS, DDoS, fuerza bruta, ataques web, infiltración, escaneo de puertos y botnet. Proporciona más de 80 características de flujo extraídas mediante su herramienta, CICFlowMeter \cite{sharafaldin2018cicids2017}.

Un año más tarde surge CSE-CIC-IDS2018, el cual tiene un mayor tamaño y está orientado a la evaluación de IDS en contextos de big data. Este dataset contiene tráfico de red recogido durante diez días y más de dieciséis millones de muestras, contando con una proporción del 17\% de tráfico de ataque aproximadamente. Su estructura es similar al CIC-IDS2017, pero fue generado a partir de una infraestructura más grande \cite{leevy2020csecicids2018}.

La principal limitación en estos conjuntos de datos más modernos es el desbalanceo entre clases, debido a que las muestras pertenecientes al tráfico benigno suelen ser más abundantes que las muestras pertenecientes a las clases minoritarias. Esto puede causar que se obtenga una precisión global elevada mientras que las clases minoritarias están siendo clasificadas incorrectamente. Por ejemplo, en el caso del CSE-CIC-IDS2018, se reportan resultados globales muy altos, pero que podrían estar afectados por sobreajuste, desbalanceo o una insuficiente limpieza de los datos. Esto indica la necesidad de un correcto preprocesamiento y división de los datos para poder interpretar correctamente los resultados \cite{leevy2020csecicids2018}.
